 $S=\set{(m,n) \in \mathbb{N}^2 \mid n!! + f(n) = m!}$とおく. 正の整数$n$を固定したとき, $(m,n)$が方程式を満たすような$m$は高々1個であるから, 十分大きい$n$から$m$が存在しないことを示せば十分である. まず, $f(x)$が零多項式である場合に主張を示す.

\begin{claim}
$n!! = m!$を満たす$(m,n)$は有限個である($(1,1), (2,2)$に限る).
\end{claim}
\begin{proof}
 $m\geq 3$の範囲に解がないことを示せばよい. このとき$n$は偶数である. $n!! = m!$より$n! > m!$なので$n>m$.
 $m!$が偶数なので$n$は偶数でなければならない. よって$(n-1)!!$は奇数であり, $v_2(m!) = v_2(n!!) = v_2(n!!(n-1)!!) = v_2(n!)$である. しかし$n>m$かつ$n$が偶数であることから$v_2(n!) > v_2(m!)$となるはずなので矛盾である.
\end{proof}


% まず$f(x)=0$の場合を示す.

% \begin{claim}
% $f(x) = 0$のとき, $S$は有限集合.
% \end{claim}
% \begin{proof} $m\geq 3$の範囲に解がないことを示せばよい. このとき$n$は偶数である. $n!! = m!$より$n! > m!$なので$n>m$.
% $m!$が偶数なので$n$は偶数でなければならない. よって$(n-1)!!$は奇数であり, $v_2(m!) = v_2(n!!) = v_2(n!!(n-1)!!) = v_2(n!)$である. しかし$n>m$かつ$n$が偶数であることから$v_2(n!) > v_2(m!)$となるはずなので矛盾である.
% \end{proof}
% 以降$f(x)$の次数は$d\geq 0$とする. このとき, $f(x)$は$d+1$個の一次式$x-2i$ ($i=0,1,\dots, d$) のいずれかでは割り切れない. つまり, $f(2i) \neq 0$となる$i$があるので, それを固定する.
% \bolm{以降は n\geq 2d+2 かつ n\geq |f(2i)| + 2i + 1で考える.} このとき$n-2i \geq 2$である. $f(x)$は整数係数なので, $f(x) = \sum_{k=0}^{d} a_k(x-2i)^{k}$ ($a_k$は整数, $a_0 = f(2i)$) と表せるが, このとき
% \begin{equation} \label{eqn:f(n)}
% f(n) \equiv f(2i) \pmod{n-2i}
% \end{equation}
% である.
%また,
以降$f(x)$は零でない多項式とすると, 十分大きい実数$R$が存在し, $n>R$において%$f(x)$は広義単調関数かつ
$|f(x)| > 0$でかつ
$-f(n) < n!! - (n-2)!!$であるようにできる. 後者については, $\displaytyle\lim_{n\to \infty} \dfrac{-f(n)}{n!!} = 0$であるからである.
%(このとき$|f(x)|$は少なくとも広義単調増加関数である).
\bolm{以降 n > Rで考える. }
\begin{claim}
%$n$は$n\geq 2d+2$ かつ $n\geq |f(2i)| + 2i + 1$ かつ
$n>R$, $(m,n) \in S$とする. このとき$n<2m+2$である.
\end{claim}

\begin{proof}
%\begin{enumerate}
    % \item もし$m\geq n$だとすると, $m \geq n-2i$でもあるので$m!$は$n-2i$の倍数である. $n!!$もまた$n-2i >0$ の倍数であるので
    % \[
    % n!! + f(n) = m! \implies 0+ f(n) \equiv 0 \pmod{n-2i} \implies f(2i) \equiv 0\pmod{n-2i}
    % \]
    % となる. $f(2i) \neq 0$より$|f(2i)| \geq n-2i$が従うが, $n-2i \geq |f(2i)| + 1$なので矛盾.
    % \item $n>R$なので$f(n) > 0$. よって,  $n!! < n!! + f(n) = m! < (2m)!!$により得られる.
    % \item
    $(n-2)!! < n!! + f(n) = m! < (2m)!!$より$n<2m+2$を得る.
%\end{enumerate}
\end{proof}

任意の素数$p$に対して$v_p(m!) = v_{p}(n!! + f(n))$である. 以下では$p=3$として固定し($p\neq 2$ならなんでもよい), $v_3(m!)$と$v_3(n!!)$の大小で場合分けを行う.
\begin{enumerate}
\item $v_3(m!) \geq v_3(n!!)$であるとき, $v_3(f(n)) \geq v_3(n!!)$でなければならない. $n$以下の整数に$3$の倍数は$[n/3]$個あり, そのうち約半数が$n!!$の積の中に現れる. よって$v_3(n!!) \geq \dfrac{1}{2}([n/3] - 1)$なので,
\[
|f(n)| \geq 3^{v_3(f(n))} \geq 3^{v_3(n!!)} \geq 3^{\frac{1}{2}([n/3] - 1)} \geq 3^{\frac{n-6}{6}}
\]
となる(最初の不等号のために$f(n) \neq 0$を仮定した). これを満たす$n$は有限個である.
\item
$v_3(m!) < v_3(n!!)$のとき, $v_3(f(n)) = v_3(m!)$であるが,
\[
|f(n)| \geq 3^{v_3(f(n))} = 3^{v_3(m!)} \geq 3^{\frac{m}{3}}
\]
となる(最初の不等号のために$f(n) \neq 0$を仮定した). %もし$f(x)$の最高次係数が正であるなら,
$n<2m+2$なので
\[
|f(n)| \geq 3^{\frac{m}{3}} > 3^{\frac{n-2}{6}}
\]
となるが, 最右辺は指数関数の増大度を持つので, これを満たす$n$は有限個である.
\end{enumerate}
 以上より, $|f(n)| \leq 3^{\frac{n-2}{6}} ( > 3^{\frac{n-6}{6}})$を満たす$n$に対しては, $v_3(m!)$と$v_3(n!)$の大小にかかわらず矛盾が生じるので, $(m,n)\notin S$が得られる. よって, $S$は有限集合である.
